%!TEX root = main.tex
\section{Data and Lexicons}
\label{sec:data}

\subsection{Vietnamese Corpus (ViHSD)}
\label{sec:vihsd}

We use UIT-ViHSD \cite{luu2021vihsd}, a corpus of Vietnamese social-media
comments annotated with three mutually exclusive labels: \clean{} (no offensive
content), \off{} (offensive language that does not target a protected group)
and \hate{} (hateful or seriously derogatory content).  The \off{}/\hate{}
distinction depends on target and intensity rather than on a particular word,
so a purely lexical rule engine is insufficient on its own.

After cleaning, the corpus contains 33{,}398 comments.  We use the fixed split
shipped with the project, and every system---rule engine, backbone, fusion
variant and analysis script---reads the same files.  Each run directory stores
the SHA-256 digest of the split it used and re-checks it before training,
ruling out a silently re-shuffled split.  Training contains 24{,}046 comments
(2{,}556 \hate{}, 1{,}605 \off{}, 19{,}885 \clean{}), validation 2{,}672
(270/212/2{,}190) and test 6{,}680 (688/444/5{,}548).

The test split is imbalanced: 83.05\% \clean{}, 10.30\% \hate{} and 6.65\%
\off{}.  Accuracy is therefore a poor headline metric---a constant \clean{}
predictor already reaches 0.8305---so macro-F1 is primary, always reported
with per-class F1.  The 6{,}680 test rows contain 6{,}571 distinct texts,
i.e. 109 duplicated rows.  They are kept because removing them would change
the frozen test set; one appears twice among the 15 regressions, which
therefore cover 14 distinct texts.

\subsection{English Corpus (Davidson et al.)}
\label{sec:davidson}

For English we use the Hate Speech and Offensive Language dataset
\cite{davidson2017automated} (tdavidson/hate\_speech\_offensive on Hugging Face).
The original labels map 0 $\to$ \hate{}, 1 $\to$ \off{}, 2 $\to$ \clean{}.
We use the fixed split in \texttt{data/splits\_en/} produced by the
reproducible pipeline: train 19{,}816 (1{,}144 \hate{}, 15{,}342 \off{},
3{,}330 \clean{}), validation 2{,}477, test 2{,}478.  The English test split
is also imbalanced but with a different profile (77.8\% \off{}, 5.8\% \hate{},
16.4\% \clean{}).

\subsection{Mixed Evaluation Sample}
\label{sec:mixed_sample}

To evaluate the bilingual pipeline end-to-end we draw a stratified random
sample (seed 42) of 500 Vietnamese and 500 English comments from the held-out
test splits, yielding the composition in Table~\ref{tab:mixed_composition}.
This sample preserves the per-language class priors and is used exclusively
for the end-to-end experiment in Sect.~\ref{sec:e2e}.

\begin{table}[H]
\centering
\caption{Mixed evaluation sample composition ($n=9{,}158$ full test splits)}
\label{tab:mixed_composition}
\begin{tabularx}{\textwidth}{l@{\extracolsep{\fill}}ccccr}
\toprule
Language & Total & \hate & \off & \clean & Source \\
\midrule
Vietnamese & 6{,}680 & 688 & 444 & 5{,}548 & ViHSD test \\
English    & 2{,}478 & 143 & 1{,}918 & 417 & Davidson test \\
\midrule
\textbf{Total} & \textbf{9{,}158} & \textbf{831} & \textbf{2{,}362} & \textbf{5{,}965} & \\
\bottomrule
\end{tabularx}
\end{table}


\subsection{Vietnamese Hate-Speech Lexicon}
\label{sec:lexicon_vi}

The Vietnamese rule engine is driven by a hate-speech lexicon of \textbf{668
entries}, mined and labelled from the \emph{training} portions of three
corpora: ViHSD; R-ViHSD (a cleaned release adding TDTU YouTube and voz-forum
comments) \cite{bontoriTdtu}; and the code-mixed ViCM corpus
\cite{nguyen2026vicm}.  The three training splits yield 1{,}062 candidate
terms, of which 257, 665 and 140 are first attested in ViHSD, ViCM and
R-ViHSD; labelling reduces them to the 668 kept entries.  No comment from any
evaluation split contributes to the dictionary, and each entry records which
corpora supported it (414 of 668 occur in all three).

Each entry stores the surface term and a canonical form, a lexical label
(\hate{} or \off{}) with a per-label severity, an entry confidence estimated
from the training data, a context-required flag, and a review status
(\emph{verified} or \emph{auto}); both statuses may match, but the flag is
retained in the evidence record.

The lexicon is deliberately not a list of slurs.  Most entries---551 of 668
(82.5\%)---are multi-word constructions such as \emph{mả cha}, \emph{con mẹ}
or \emph{ngu hơn}, and 524 (78.4\%) are context-required, because Vietnamese
insults frequently work by comparison and implication rather than by a single
taboo word.  Only 74 entries carry an \off{} lexical label against 594 for
\hate{}, which is why the rule engine is weak on \off{} and fusion keeps that
channel mostly neural.

\subsection{English Hate-Speech Lexicon}
\label{sec:lexicon_en}

The English rule engine uses a lexicon of \textbf{818 entries}, built by merging
a 40-term manual seed set with 778 entries extracted from the Davidson training
split (19{,}816 documents) using relaxed statistical parameters (min.\ doc
frequency 5, min.\ lift 1.2, min.\ HATE precision 0.20, min.\ OFFENSIVE precision
0.60, min.\ Wilson lower bound 0.15, min.\ prior margin 0.02).  The extraction
result, stored in \texttt{verified\_lexicon\_en\_combined.csv} (version
\texttt{v1-en-combined}) and configured via \texttt{configs/rule\_engine\_en.yaml},
contains 86 \hate{}, 356 \off{}, and 376 \clean{} terms.

The English lexicon is larger in raw count (818 vs. 668) but contains far fewer
HATE terms (86 vs. 594) because the Davidson corpus has only 5.8\% HATE
documents.  This asymmetry directly impacts the hybrid's HATE F1 (0.46 EN vs
0.61 VI on the mixed sample) and is the highest-priority target for future
lexicon expansion.

\paragraph{Ethics of examples.}
All quoted comments are public test comments.  Because they contain profanity
and slurs, we mask offending characters with \texttt{*} in every table;
evaluation always runs on the unmasked text.  No user identifiers are
reproduced, and names that occasionally appear inside comments are part of the
public corpus text.